\documentclass[11pt,oneside]{article}
\usepackage[a4paper,margin=25mm,headheight=15pt]{geometry}
\usepackage[T1]{fontenc}
\usepackage[utf8]{inputenc}
\IfFileExists{libertinus.sty}{\usepackage{libertinus}}{\usepackage{lmodern}}
\usepackage{microtype}
\usepackage{amsmath,amssymb,amsthm,mathtools,mathrsfs,bm}
\usepackage{booktabs,longtable,tabularx,array}
\usepackage{xcolor}
\usepackage{enumitem}
\usepackage{fancyhdr}
\usepackage{titlesec}
\usepackage{xurl}
\usepackage{hyperref}
\usepackage[nameinlink,noabbrev]{cleveref}
\usepackage{listings}

\definecolor{linkblue}{RGB}{25,75,135}
\definecolor{shadegray}{RGB}{246,247,249}
\definecolor{rulegray}{RGB}{195,201,208}
\definecolor{provedgreen}{RGB}{28,112,74}
\definecolor{conjpurple}{RGB}{105,55,135}
\hypersetup{
  colorlinks=true,
  linkcolor=linkblue,
  citecolor=linkblue,
  urlcolor=linkblue,
  pdftitle={Automatic Spectra, Exact Dyadic Cubature, and Probabilistic Duals in the Pascal--Rvachev Hierarchy},
  pdfsubject={Repository-audited frontier results on Fabius and Rvachev functions},
  pdfkeywords={Fabius function, Rvachev up-function, Thue-Morse, Mahler function, hypertranscendence, natural boundary, Poisson-binomial, Laguerre-Polya, q-series}
}
\pagestyle{fancy}
\fancyhf{}
\fancyhead[L]{\small Pascal--Rvachev frontiers}
\fancyhead[R]{\small\nouppercase{\leftmark}}
\fancyfoot[C]{\thepage}
% Canonical project-wide mathematical notation.
% Canonical mathematical notation for active FabiusFunction LaTeX documents.
%
% This file is intentionally a plain TeX input rather than a LaTeX package.
% Every standalone document loads it by an explicit path relative to that
% document's directory; included fragments inherit the definitions from their
% root document.  Keep this file independent of document layout and metadata.
%
% Contract:
%   * one command has one mathematical meaning throughout the active corpus;
%   * names encode semantic roles rather than a document-local choice of letter;
%   * visually similar but differently normalized objects have different print;
%   * every command below is defined normatively in the notation catalogue.
%
% The active corpus excludes docs/archive/.  Historical sources may retain
% their original notation only inside an explicitly labelled quotation.

\makeatletter
\@ifundefined{mathbb}{%
  \PackageError{fabius-notation}{The amssymb package must be loaded first}%
    {Load amsmath and amssymb before inputting fabius-notation.tex.}%
}{}
\@ifundefined{operatorname}{%
  \PackageError{fabius-notation}{The amsmath package must be loaded first}%
    {Load amsmath before inputting fabius-notation.tex.}%
}{}
\makeatother

% Version stamp used by the catalogue and by static audits.
\newcommand{\FabiusNotationVersion}{2026-08-31}

% Number systems.  NaturalNumbers is explicitly the nonnegative convention.
\newcommand{\NaturalNumbers}{\mathbb{N}_{0}}
\newcommand{\PositiveIntegers}{\mathbb{N}_{>0}}
\newcommand{\IntegerNumbers}{\mathbb{Z}}
\newcommand{\RationalNumbers}{\mathbb{Q}}
\newcommand{\RealNumbers}{\mathbb{R}}
\newcommand{\ComplexNumbers}{\mathbb{C}}
\newcommand{\ExtendedRealNumbers}{\overline{\mathbb{R}}}

% The three principal Fabius--Rvachev functions.  Subscripts are deliberate:
% a bare F and a calligraphic F have several unrelated uses in the corpus.
\newcommand{\FabiusBounded}{F_{\mathrm{bd}}}
\newcommand{\FabiusGlobal}{F_{\mathrm{gl}}}
\newcommand{\FabiusQuantile}{Q_{F}}
\newcommand{\FabiusClampedQuantile}{Q_{F,\mathrm{cl}}}
\newcommand{\RvachevUp}{\operatorname{up}}

% Constants, elementary operators, and delimiters.
\newcommand{\EulerE}{\mathrm{e}}
\newcommand{\ImaginaryUnit}{\mathrm{i}}
\newcommand{\Differential}{\mathop{}\!\mathrm{d}}
\newcommand{\AbsoluteValue}[1]{\left\lvert #1\right\rvert}
\newcommand{\Norm}[1]{\left\lVert #1\right\rVert}
\newcommand{\Floor}[1]{\left\lfloor #1\right\rfloor}
\newcommand{\Ceiling}[1]{\left\lceil #1\right\rceil}
\newcommand{\FractionalPart}[1]{\left\{#1\right\}}
\newcommand{\PositivePart}[1]{\left(#1\right)_{+}}
\newcommand{\IndicatorOf}[1]{\mathbf{1}_{#1}}
\newcommand{\SetOf}[1]{\left\{#1\right\}}
\newcommand{\InnerProduct}[1]{\left\langle #1\right\rangle}
\newcommand{\CardinalityOf}[1]{\left\lvert #1\right\rvert}
\newcommand{\DefinitionEquals}{\mathrel{\coloneqq}}
\newcommand{\Sign}{\operatorname{sgn}}
\newcommand{\IdentityOperator}{\operatorname{id}}
\newcommand{\SpanOperator}{\operatorname{span}}
\newcommand{\TraceOperator}{\operatorname{tr}}
\newcommand{\RankOperator}{\operatorname{rank}}
\newcommand{\DiagonalOperator}{\operatorname{diag}}
\newcommand{\ResidueOperator}{\operatorname*{Res}}
\newcommand{\LeastCommonMultiple}{\operatorname{lcm}}
\newcommand{\OrderOperator}{\operatorname{ord}}
\newcommand{\PrincipalLogarithm}{\operatorname{Log}}
\newcommand{\FallingFactorial}[2]{\left(#1\right)^{\underline{#2}}}
\newcommand{\RisingFactorial}[2]{\left(#1\right)^{\overline{#2}}}

% Support, topology, convolution, and metric notation.
\newcommand{\SupportOperator}{\operatorname{supp}}
\newcommand{\SupportOf}[1]{\SupportOperator\!\left(#1\right)}
\newcommand{\TopologicalSupportOf}[1]{\operatorname{tsupp}\!\left(#1\right)}
\newcommand{\Convolution}{\mathbin{*}}
\newcommand{\MetricDistance}{\operatorname{dist}}
\newcommand{\TotalVariationDistance}{d_{\mathrm{TV}}}

% Probability.  Equality and convergence in law are not metric distance.
\newcommand{\Probability}{\mathbb{P}}
\newcommand{\Expectation}{\mathbb{E}}
\newcommand{\Variance}{\operatorname{Var}}
\newcommand{\Covariance}{\operatorname{Cov}}
\newcommand{\LawOperator}{\operatorname{Law}}
\newcommand{\LawOf}[1]{\LawOperator\!\left(#1\right)}
\newcommand{\EqualInLaw}{\mathrel{\overset{\mathrm{law}}{=}}}
\newcommand{\ConvergesInLaw}{\mathrel{\xrightarrow{\mathrm{law}}}}

% Fourier and integral transforms.  The printed labels expose normalization.
% FourierTwoPi uses exp(-2*pi*i*x*xi); FourierAngular uses exp(-i*x*omega).
\newcommand{\FourierTwoPi}[1]{\widehat{#1}_{\,2\pi}}
\newcommand{\FourierAngular}[1]{\widehat{#1}_{\,\mathrm{ang}}}
\newcommand{\LaplaceTransformSymbol}{\mathcal{L}}
\newcommand{\MellinTransformSymbol}{\mathcal{M}}
\newcommand{\LaplaceTransformOf}[1]{\LaplaceTransformSymbol\!\left[#1\right]}
\newcommand{\MellinTransformOf}[1]{\MellinTransformSymbol\!\left[#1\right]}
\newcommand{\MomentGeneratingFunction}[1]{M_{#1}}
\newcommand{\CharacteristicFunction}[1]{\varphi_{#1}}

% The two standard sinc functions must never share one symbol.
\newcommand{\SincRad}{\operatorname{sinc}_{\mathrm{rad}}}
\newcommand{\SincPi}{\operatorname{sinc}_{\pi}}

% Binary arithmetic and Thue--Morse notation.
\newcommand{\BinaryDigitSum}{\operatorname{wt}_{2}}
\newcommand{\ThueMorseSignSymbol}{\varepsilon}
\newcommand{\ThueMorseSign}[1]{\varepsilon_{#1}}
\newcommand{\PadicValuation}[1]{\nu_{#1}}
\newcommand{\TwoAdicValuation}{\nu_{2}}

% q-series notation.  Every base is an explicit argument.
\newcommand{\QPochhammer}[3]{\left(#1;#2\right)_{#3}}
\newcommand{\GaussianBinomial}[3]{%
  \genfrac{[}{]}{0pt}{}{#1}{#2}_{#3}%
}
\newcommand{\QInteger}[2]{\left[#1\right]_{#2}}

% Named special functions.
\newcommand{\Polylogarithm}{\operatorname{Li}}
\newcommand{\LambertW}{W}
\newcommand{\GammaFunction}{\Gamma}
\newcommand{\RiemannZeta}{\zeta}

% Asymptotics.  BigO and LittleO are operator tokens for existing formula
% grammar; prefer the At forms so that the limiting regime is printed.
\newcommand{\BigO}{\mathcal{O}}
\newcommand{\LittleO}{o}
\newcommand{\BigOAt}[2]{\mathcal{O}_{#2}\!\left(#1\right)}
\newcommand{\LittleOAt}[2]{o_{#2}\!\left(#1\right)}

\renewcommand{\headrulewidth}{0.4pt}
\titleformat{\section}{\Large\bfseries}{\thesection}{0.7em}{}
\titleformat{\subsection}{\large\bfseries}{\thesubsection}{0.7em}{}
\titleformat{\subsubsection}{\normalsize\bfseries}{\thesubsubsection}{0.7em}{}
\setcounter{secnumdepth}{3}
\setcounter{tocdepth}{2}
\setlist[itemize]{leftmargin=2em,itemsep=0.25em,topsep=0.4em}
\setlist[enumerate]{leftmargin=2.4em,itemsep=0.25em,topsep=0.4em}
\setlength{\emergencystretch}{3em}

\newtheorem{theorem}{Theorem}[section]
\newtheorem{proposition}[theorem]{Proposition}
\newtheorem{lemma}[theorem]{Lemma}
\newtheorem{corollary}[theorem]{Corollary}
\newtheorem{conjecture}[theorem]{Conjecture}
\theoremstyle{definition}
\newtheorem{definition}[theorem]{Definition}
\newtheorem{example}[theorem]{Example}
\newtheorem{problem}[theorem]{Problem}
\theoremstyle{remark}
\newtheorem{remark}[theorem]{Remark}
\newtheorem{warning}[theorem]{Warning}

\newcommand{\Po}{\operatorname{Po}}
\newcommand{\Normal}{\mathcal N}
\newcommand{\LP}{\mathcal{LP}}
\newcommand{\one}{\mathbf 1}
\newcommand{\ZetaR}{\mathcal Z}
\newcommand{\Ar}{\mathcal A}
\newcommand{\Mr}{\mathcal M}
\newcommand{\Bell}{\mathcal B}
\newcommand{\New}{\textcolor{provedgreen}{\textsf{new relative to audited repository}}}
\newcommand{\Conj}{\textcolor{conjpurple}{\textsf{conjectural}}}
\newcommand{\Known}{\textsf{repository baseline}}

\newenvironment{resultbox}{%
  \par\smallskip\noindent\begin{center}\begin{minipage}{0.94\textwidth}%
  \hrule\vspace{0.65em}\small
}{%
  \vspace{0.65em}\hrule\end{minipage}\end{center}\smallskip
}

\lstdefinestyle{code}{
  basicstyle=\ttfamily\small,
  backgroundcolor=\color{shadegray},
  frame=single,
  rulecolor=\color{rulegray},
  columns=fullflexible,
  keepspaces=true,
  showstringspaces=false,
  xleftmargin=0.6em,
  xrightmargin=0.6em
}

\title{\bfseries Automatic Spectra, Exact Dyadic Cubature,\\
 and Probabilistic Duals in the Pascal--Rvachev Hierarchy\\[0.5em]
\large A repository-audited frontier report on Fabius, Rvachev, Thue--Morse,\\
Mahler, $q$-series, Bell--Bernoulli, and spectral-product structures}
\author{Prepared from the \texttt{ProveIt} Fabius-function document corpus}
\date{30 August 2026}

\begin{document}
\maketitle

\begin{abstract}
This report audits the current LaTeX research corpus under
\path{Analysis/FabiusFunction/docs} in Vladimir Reshetnikov's
\texttt{ProveIt} repository and develops a set of results not located in the
consolidated documents or their manifest descriptions.  The audit treats the
primary 9509-line exposition, the 14587-line canonical frontier volume, the
spectra-and-arithmetic, exponent-and-$q$-series, dyadic-comb, and frontier-compilation
volumes, and the repository's draft manifest as the principal current
representatives of the recursively organized corpus.  Historical drafts that
have been editorially absorbed are tracked through that manifest.  The word
``new'' below therefore means \emph{new relative to this audited repository
state}; it is not an unconditional claim of first publication in the world
literature.

The first main result closes a problem explicitly left open in the repository:
for every rank $r\geq1$, the generalized Thue--Morse spectral-sign generating
function $G_r$ is nonrational, hypertranscendental over $\ComplexNumbers(z)$, and has the
unit circle as a natural boundary.  The proof combines a new nonperiodicity
argument for strongly block-multiplicative digit characters with the
Adamczewski--Dreyfus--Hardouin Mahler dichotomy.  A parallel construction for
the valuation-multiplicity Lambert series
\[
 L_r(q)=\sum_{n\geq1}\binom{\TwoAdicValuation(n)+1}{r}q^n
\]
yields an explicit order-$(r+1)$ homogeneous Mahler equation, a dense set of
radial singularities at dyadic roots of unity, hypertranscendence, and
non-$D$-finiteness.  A bivariate master Lambert germ packages all ranks.

The analytic zero divisor then gives a sharp higher-rank Strang--Fix theorem:
on the lattice $Q^{-1}\IntegerNumbers$, the density $f_r$ integrates every polynomial of
degree strictly below $\binom{\TwoAdicValuation(Q)+1}{r}$ by an exact finite periodized
comb, and the threshold is optimal.  Appell reconstruction converts this into
exact polynomial reproduction, with immediate Legendre and Bernoulli
corollaries.

Wick rotation of the canonical product produces
\[
 \Ar_r(w)=\Phi_r(\ImaginaryUnit\sqrt w)
 =\prod_{n\geq1}\left(1+\frac{w}{n^2}\right)^{\mu_r(n)}.
\]
For every rank this is a type-I Laguerre--P\'olya function of order $1/2$;
its coefficient sequence is $PF_\infty$, so all Toeplitz minors, shifted
Jensen hyperbolicity statements, and ultra-log-concavity inequalities lift
from the known rank-one case to the full Pascal hierarchy.

A new probability interpretation is especially productive.  The normalized
ratio $\Ar_r(\theta z)/\Ar_r(\theta)$ is the probability generating function
of an infinite Poisson-binomial spectral count $N_{r,\theta}$.  It identifies
exact even moments with count probabilities, obeys a rank-branching law, and
has a multi-scale limit theory.  For $\theta_r=c a^r$, the count is degenerate
when $a<3$, converges in total variation to a fixed Poisson law at $a=3$, is
asymptotically Poisson with diverging mean for $3<a<\sqrt{15}$, has a
nontrivial mod-Poisson correction at $a=\sqrt{15}$, and is asymptotically
Gaussian throughout $3<a<5$.  Finally, $\Ar_r(w)^{-\tau}$ is the Laplace
transform of a generalized-gamma convolution with an explicit Thorin measure;
a separate rank central limit theorem and geometric standardized-cumulant
decay follow.

The report supplies complete proofs, exact Bell--Bernoulli formulas,
computational checks, a reproducible Python program, a conjecture register,
and a staged formalization plan.
\end{abstract}

\tableofcontents
\newpage

\section{Scope, audit method, and originality boundary}

\subsection{What was audited}

The current repository is not a small collection of independent notes.  Its
own manifest states that absorbed drafts are deleted after their content has
been integrated into a primary or consolidated frontier volume, with the
manifest and git history serving as the inventory and archive
\cite{repo-manifest}.  Accordingly, the audit used two complementary layers.

\begin{enumerate}
\item The primary human-readable exposition
\path{Fabius_Function_and_Rvachev_Up.tex}, which separates the bounded Fabius
CDF, Rvachev's compactly supported bump, and the signed global extension, and
collects the proved probability, Fourier, arithmetic, inverse, Legendre, and
Lambert-asymptotic results \cite{repo-primary}.
\item The current consolidated frontier documents, especially
\path{semi-formalized-research-frontiers.tex},
\path{Spectra_and_Arithmetic_Frontiers.tex},
\path{Exponents_and_q_Series_Frontiers.tex},
\path{Dyadic_Comb_Frontiers.tex}, and
\path{Frontier_Compilations.tex} \cite{repo-frontier,repo-spectra,repo-q,repo-comb,repo-compilations}.
The manifest was then used to identify topics already absorbed from historical
or incoming drafts.
\end{enumerate}

This method is stronger than searching a few obvious phrases and weaker than a
formal semantic comparison of every theorem in every historical commit.  The
novelty labels in this report are therefore deliberately local: they record
absence from the audited current repository corpus, not an exhaustive global
priority search.

\subsection{Important duplication exclusions}

Several attractive conclusions were initially rediscovered and then removed
from the novelty ledger because the corpus already contains them.  In
particular:

\begin{itemize}
\item the canonical product for the classical Rvachev transform, its integer
zero multiplicities $1+\TwoAdicValuation(n)$, and exact zero jets;
\item the Pascal--Rvachev hierarchy, its random-series model, support,
refinement law, multiplicity formula
$\mu_r(n)=\binom{\TwoAdicValuation(n)+1}{r}$, spectral zeta function, digital sign law,
and Mahler equation \cite{repo-spectra};
\item rank-one sharp dyadic-comb moment exactness and Appell reproduction
\cite{repo-comb,repo-compilations};
\item the rank-one type-I Laguerre--P\'olya classification, $PF_\infty$
coefficient theorem, Jensen hyperbolicity, moment inequalities,
non-infinite-divisibility result, and gamma dual \cite{repo-compilations};
\item higher-rank zero counting, spectral zeta, heat-trace, endpoint, and
Fourier-decay results already developed in the spectra volume.
\end{itemize}

The present results use those statements as inputs.  They do not rebrand them
as new.

\subsection{Result ledger}

\begin{longtable}{@{}p{0.26\textwidth}p{0.18\textwidth}p{0.48\textwidth}@{}}
\toprule
Topic & Status & Contribution in this report\\
\midrule
\endhead
Spectral signs $G_r$ & \New & Proof of non-eventual-periodicity, nonrationality, hypertranscendence, and the unit-circle natural boundary; this resolves Problem D(iv) stated in the spectra volume.\\
Multiplicity Lambert series $L_r$ & \New & Divisibility-layer formula, rank recurrence, order-$(r+1)$ Mahler equation, dense radial singularities, natural boundary, hypertranscendence, and non-$D$-finiteness.\\
Rank master Lambert germ & \New & The bivariate identity $\mathscr L(u,q)=\sum_h(1+u)^h q^{2^h}/(1-q^{2^h})$ and its first-order Mahler equation.\\
Higher-rank comb exactness & \New & Sharp degree $\binom{\TwoAdicValuation(Q)+1}{r}-1$ for polynomial cubature and Appell reconstruction on $Q^{-1}\IntegerNumbers$.\\
Laguerre--P\'olya/PF hierarchy & \New & Lift of the rank-one product theorem to every Pascal rank, including all Toeplitz-minor and moment inequalities.\\
Poisson-binomial spectral count & \New & Exact PGF, moment-probability duality, factorial cumulants, rank branching, saddle tilting, Poisson/mod-Poisson/Gaussian phase diagram.\\
Generalized-gamma dual & \New & All-rank Thorin measure, rank branching, $L^2$ concentration, and a standardized rank CLT.\\
Hypertranscendence of $\Phi_r$ or $\Ar_r$ & \Conj & Non-$D$-finiteness is proved; full differential transcendence remains conjectural.\\
Exact higher-rank dyadic values and inverse asymptotics & Open & Concrete directions are given, but no theorem is claimed here.\\
\bottomrule
\end{longtable}

\begin{resultbox}
\textbf{Proof-status convention.}
A theorem or corollary below is proved in the report.  A numerical table is
supporting evidence only.  Every conjecture is marked explicitly.  The
Python program verifies finite identities and numerical predictions but is
not used as a substitute for any analytic proof.
\end{resultbox}

\section{Baseline: the Pascal--Rvachev hierarchy}

\subsection{Fourier convention and random series}

Throughout this report
\[
 \SincPi z:=\frac{\sin(\pi z)}{\pi z},\qquad \SincPi0=1.
\]
The normalized convention puts the zeros at the nonzero integers.

\begin{definition}[Pascal--Rvachev products]
Set $\Phi_0(z)=\SincPi(2z)$.  For $r\geq1$, let
\begin{equation}\label{eq:Phi-r}
 a_{r,h}:=\binom{h}{r-1},\qquad
 \Phi_r(z):=\prod_{h\geq0}\SincPi\!\left(\frac{z}{2^h}\right)^{a_{r,h}}.
\end{equation}
The convention $\binom hk=0$ for $h<k$ is always in force.
\end{definition}

The repository proves local uniform convergence and gives the random-series
model
\begin{equation}\label{eq:X-r}
 X_r=\sum_{h\geq0}\sum_{\ell=1}^{a_{r,h}}U_{h,\ell},
 \qquad
 U_{h,\ell}\sim\operatorname{Unif}[-2^{-h-1},2^{-h-1}],
\end{equation}
with all summands independent, so that
\begin{equation}\label{eq:characteristic}
 \Expectation\EulerE^{-2\pi\ImaginaryUnit zX_r}=\Phi_r(z),\qquad \SupportOperator X_r=[-1,1].
\end{equation}
Write $f_r$ for the density.  The case $r=1$ is the classical Rvachev
up-density in the Fourier normalization used here; its translate/integral is
the bounded Fabius function.

Pascal's identity gives the refinement ladder
\begin{equation}\label{eq:refinement}
 \Phi_r(2z)=\Phi_r(z)\Phi_{r-1}(z),
 \qquad
 2X_r\stackrel d=X_r+X_{r-1}',
\end{equation}
where the prime denotes an independent copy.

\subsection{Zero divisor and spectral zeta}

For $n\neq0$, the zero order at $n$ is
\begin{equation}\label{eq:mu-r}
 \mu_r(n)=\sum_{h=0}^{\TwoAdicValuation(|n|)}\binom h{r-1}
 =\binom{\TwoAdicValuation(|n|)+1}{r}.
\end{equation}
Consequently
\begin{equation}\label{eq:canonical}
 \Phi_r(z)=\prod_{n\geq1}
 \left(1-\frac{z^2}{n^2}\right)^{\mu_r(n)}.
\end{equation}
The associated Dirichlet series is
\begin{equation}\label{eq:spectral-zeta}
 \ZetaR_r(s):=\sum_{n\geq1}\frac{\mu_r(n)}{n^s}
 =\zeta(s)\frac{2^{-(r-1)s}}{(1-2^{-s})^r}
 =\frac{2^s\zeta(s)}{(2^s-1)^r},\qquad \Re s>1.
\end{equation}
These are repository baselines \cite{repo-spectra}.

Two elementary consequences will be used repeatedly:
\begin{equation}\label{eq:zeta-special}
 \ZetaR_r(2)=\frac{4\zeta(2)}{3^r},\qquad
 \ZetaR_r(4)=\frac{16\zeta(4)}{15^r},\qquad
 \ZetaR_r(2k)=\frac{2^{2k}\zeta(2k)}{(2^{2k}-1)^r}.
\end{equation}
Also, the first positive zero is $2^{r-1}$, and
\begin{equation}\label{eq:unbounded-multiplicity}
 \mu_r(2^H)=\binom{H+1}{r}\longrightarrow\infty.
\end{equation}

\subsection{Moments, cumulants, and Bell polynomials}

Let $m_{r,n}=\Expectation X_r^n$.  Symmetry kills odd moments.  Taking logarithms in
\eqref{eq:canonical} gives, for $|z|<2^{r-1}$,
\begin{equation}\label{eq:log-Phi}
 \log\Phi_r(z)=-\sum_{k\geq1}\frac{\ZetaR_r(2k)}{k}z^{2k}.
\end{equation}
Equivalently, the even cumulants are
\begin{equation}\label{eq:cumulants-X}
 \kappa_{2k}(X_r)=(-1)^{k+1}
 \frac{2(2k-1)!}{(2\pi)^{2k}}\ZetaR_r(2k).
\end{equation}
This is another repository baseline.  We record the Bell-polynomial form
because it will connect the analytic and probabilistic results below.

Define
\begin{equation}\label{eq:A-def-early}
 \Ar_r(w):=\Phi_r(\ImaginaryUnit\sqrt w)=\sum_{n\geq0}C_{r,n}w^n.
\end{equation}
Then
\begin{equation}\label{eq:Bell-coeff}
 C_{r,n}=\frac{1}{n!}\Bell_n(x_{r,1},\ldots,x_{r,n}),
 \qquad
 x_{r,k}=(-1)^{k+1}(k-1)!\ZetaR_r(2k),
\end{equation}
where $\Bell_n$ is the complete exponential Bell polynomial, and
\begin{equation}\label{eq:C-moment}
 C_{r,n}=\frac{(2\pi)^{2n}}{(2n)!}m_{r,2n}.
\end{equation}
Euler's formula
\[
 \zeta(2k)=(-1)^{k+1}\frac{B_{2k}(2\pi)^{2k}}{2(2k)!}
\]
makes every $m_{r,2n}$ rational and expresses it explicitly through Bell
polynomials in Bernoulli numbers.

\section{Hypertranscendence and natural boundaries of the spectral signs}

\subsection{The generalized Thue--Morse characters}

Let $b_j(N)\in\{0,1\}$ be the binary digits of $N$.  The spectra volume
defines
\begin{equation}\label{eq:epsilon}
 \varepsilon_r(N)=(-1)^{\sum_{j\geq0}b_j(N)\binom j{r-1}},
 \qquad
 \sigma_{r,j}:=(-1)^{\binom j{r-1}}.
\end{equation}
Let
\begin{equation}\label{eq:Pr}
 P_r:=2^{\lceil\log_2r\rceil},\qquad B_r:=2^{P_r}.
\end{equation}
Lucas's theorem implies that $j\mapsto\sigma_{r,j}$ is periodic with period
$P_r$, and at least one entry per period is $-1$.  The generating function is
\begin{equation}\label{eq:G-product}
 G_r(z):=\sum_{N\geq0}\varepsilon_r(N)z^N
 =\prod_{j\geq0}(1+\sigma_{r,j}z^{2^j}).
\end{equation}
Splitting off one period gives the repository's Mahler equation
\begin{equation}\label{eq:G-Mahler}
 G_r(z)=A_r^{\rm sign}(z)G_r(z^{B_r}),
 \qquad
 A_r^{\rm sign}(z)=\prod_{j=0}^{P_r-1}(1+\sigma_{r,j}z^{2^j}).
\end{equation}
The open frontier list asks for the natural-boundary behavior of $G_r$ on
$|z|=1$ \cite{repo-spectra}.

\subsection{A nonperiodicity lemma for periodic digit characters}

The key new input is not a numerical observation but an elementary rigidity
argument.

\begin{lemma}[Strong block multiplicativity]\label{lem:block}
For $0\leq a<B_r^m$ and $k\geq0$,
\begin{equation}\label{eq:block}
 \varepsilon_r(a+B_r^mk)=\varepsilon_r(a)\varepsilon_r(k).
\end{equation}
\end{lemma}

\begin{proof}
The binary digits of $a$ occupy positions below $mP_r$, while those of
$B_r^mk$ are the digits of $k$ shifted by $mP_r$.  Periodicity of
$\sigma_{r,j}$ with period $P_r$ leaves every shifted digit weight unchanged.
The two sets of occupied positions are disjoint, so the signs multiply.
\end{proof}

\begin{lemma}[Eventual periodicity would be pure periodicity]\label{lem:eventual-pure}
If $\varepsilon_r$ is eventually periodic with period $d\geq1$, then it is
periodic with period $d$ from $N=0$ onward.
\end{lemma}

\begin{proof}
Suppose $\varepsilon_r(N+d)=\varepsilon_r(N)$ for $N\geq N_0$.  Choose $m$
so large that $B_r^m>N_0+d$, and then choose $k$ with $B_r^mk\geq N_0$.
For $0\leq a<B_r^m-d$, eventual periodicity and
\cref{lem:block} give
\[
 \varepsilon_r(k)\varepsilon_r(a+d)
 =\varepsilon_r(B_r^mk+a+d)
 =\varepsilon_r(B_r^mk+a)
 =\varepsilon_r(k)\varepsilon_r(a).
\]
Cancel the sign $\varepsilon_r(k)$.  Since $m$ can be arbitrarily large, the
identity $\varepsilon_r(a+d)=\varepsilon_r(a)$ holds for every $a\geq0$.
\end{proof}

\begin{theorem}[Nonperiodicity of every spectral sign sequence]\label{thm:nonperiodic}
For every $r\geq1$, the sequence $(\varepsilon_r(N))_{N\geq0}$ is not
eventually periodic.
\end{theorem}

\begin{proof}
Assume it is eventually periodic.  By \cref{lem:eventual-pure}, it is purely
periodic with some period $d$.  Fix $j$ so large that $2^j>d$.  Put
\[
 S_j:=\prod_{k=0}^{j-1}\sigma_{r,k},\qquad c:=\varepsilon_r(d-1)\in\{\pm1\}.
\]
Pure periodicity gives
\[
 \sigma_{r,j}=\varepsilon_r(2^j)=\varepsilon_r(2^j-d).
\]
In the first $j$ binary positions, $2^j-d=(2^j-1)-(d-1)$ is the bitwise
complement of $d-1$.  Hence
\begin{equation}\label{eq:complement-sign}
 \varepsilon_r(2^j-d)
 =\prod_{k=0}^{j-1}\sigma_{r,k}^{1-b_k(d-1)}
 =\frac{S_j}{c}.
\end{equation}
Therefore $\sigma_{r,j}=S_j/c$, and
\[
 S_{j+1}=S_j\sigma_{r,j}=\frac{S_j^2}{c}=\frac1c=c.
\]
For all sufficiently large $j$, both $S_j$ and $S_{j+1}$ equal $c$, so
$\sigma_{r,j}=1$ eventually.  This contradicts the periodic mask's having at
least one negative entry in every period.
\end{proof}

\subsection{Mahler hypertranscendence}

We use the following deep dichotomy of Adamczewski, Dreyfus, and Hardouin
\cite{adh}.

\begin{theorem}[Mahler differential dichotomy, cited form]\label{thm:ADH}
Let $p\geq2$, and let a possibly ramified Laurent series solve a nontrivial
linear Mahler equation over $\ComplexNumbers(z)$ involving $f(z),f(z^p),\ldots,f(z^{p^m})$.
Then either the solution belongs to the ramified rational field
$\bigcup_{d\geq1}\ComplexNumbers(z^{1/d})$, or it is hypertranscendental over $\ComplexNumbers(z)$.
\end{theorem}

A power series in integral powers that lies in a ramified rational field is in
fact rational in $z$: invariance under all deck transformations of a common
ramification forces the rational function to descend to $\ComplexNumbers(z)$.

\begin{lemma}[Finite-alphabet rational series]\label{lem:finite-rational}
A rational power series whose coefficients lie in a finite set is eventually
periodic.
\end{lemma}

\begin{proof}
Its coefficients satisfy a constant-coefficient linear recurrence from some
index onward.  A state consists of a fixed finite window of consecutive
coefficients.  Only finitely many states occur, and deterministic recurrence
propagation makes the state sequence eventually periodic.  Therefore the
coefficient sequence is eventually periodic.
\end{proof}

\begin{theorem}[Hypertranscendence and natural boundary of $G_r$]\label{thm:G-main}
For every rank $r\geq1$:
\begin{enumerate}[label=(\roman*)]
\item $G_r(z)$ is not rational;
\item $G_r(z)$ is hypertranscendental over $\ComplexNumbers(z)$;
\item the unit circle is a natural boundary of $G_r$.
\end{enumerate}
Thus the natural-boundary part of Problem D in the spectra-and-arithmetic
frontier volume has an affirmative answer for every rank.
\end{theorem}

\begin{proof}
The coefficients of $G_r$ lie in $\{\pm1\}$.  If $G_r$ were rational,
\cref{lem:finite-rational} would make $\varepsilon_r$ eventually periodic,
contradicting \cref{thm:nonperiodic}.  Equation \eqref{eq:G-Mahler} is a
nontrivial linear Mahler equation, so \cref{thm:ADH} and the descent observation
imply hypertranscendence.

For the natural boundary there are two independent routes.  First,
$G_r$ has integer coefficients and radius of convergence exactly $1$; the
P\'olya--Carlson theorem says that such a series is either rational or has the
unit circle as a natural boundary.  Second, the Mahler-specific theorem of
Bell, Coons, and Rowland gives the same rational/natural-boundary dichotomy
for a Mahler function meromorphic in the unit disk \cite{bcr}.  Nonrationality
rules out the first alternative in either theorem.
\end{proof}

\begin{corollary}[Generic algebraic independence of values]\label{cor:G-values}
Fix $m\geq0$ and a compact subset $K$ of the open unit disk.  For every
algebraic point $\alpha\in K$ outside a finite exceptional set, the numbers
\[
 G_r(\alpha),G_r'(\alpha),\ldots,G_r^{(m)}(\alpha)
\]
are algebraically independent over $\overline\RationalNumbers$.
\end{corollary}

\begin{proof}
This is the value-specialization consequence of the
Adamczewski--Dreyfus--Hardouin theorem for nonrational Mahler functions
\cite{adh}.
\end{proof}

\begin{remark}[What this does not solve]
A natural boundary is an analytic-continuation statement.  It does not by
itself determine rarefied sums, Lyapunov exponents, or the pure-point versus
singular-continuous decomposition of the diffraction measure.  The other
parts of the repository's Problem D remain open.
\end{remark}

\subsection{First masks}

\begin{center}
\begin{tabular}{@{}rrrrl@{}}
\toprule
$r$ & $P_r$ & $B_r$ & negative mask entries & $(\sigma_{r,0},\ldots,\sigma_{r,P_r-1})$\\
\midrule
1 & 1 & 2 & 1 & $-$\\
2 & 2 & 4 & 1 & $+-$\\
3 & 4 & 16 & 2 & $++--$\\
4 & 4 & 16 & 1 & $+++- $\\
5 & 8 & 256 & 4 & $++++----$\\
6 & 8 & 256 & 2 & $+++++-+-$\\
\bottomrule
\end{tabular}
\end{center}

The Python checks in the archive verify the block identity on large finite
ranges for ranks $1$ through $8$, but \cref{thm:nonperiodic} is the proof.

\section{A hypertranscendental hierarchy of valuation Lambert series}

\subsection{Divisibility layers}

Define
\begin{equation}\label{eq:L-def}
 L_r(q):=\sum_{n\geq1}\mu_r(n)q^n,
 \qquad |q|<1.
\end{equation}
The binomial multiplicity admits a useful layer-cake identity.

\begin{lemma}[Divisibility decomposition]\label{lem:mu-layer}
For $n\geq1$,
\begin{equation}\label{eq:mu-layer}
 \mu_r(n)=\sum_{h\geq r-1}\binom h{r-1}\one_{2^h\mid n}.
\end{equation}
Consequently,
\begin{equation}\label{eq:L-layer}
 \boxed{
 L_r(q)=\sum_{h\geq r-1}\binom h{r-1}
 \frac{q^{2^h}}{1-q^{2^h}}.}
\end{equation}
\end{lemma}

\begin{proof}
If $v=\TwoAdicValuation(n)$, only $h\leq v$ contribute, and the hockey-stick identity
gives
$\sum_{h=r-1}^{v}\binom h{r-1}=\binom{v+1}{r}$.  Summing
$\one_{2^h\mid n}q^n$ over $n$ gives $q^{2^h}/(1-q^{2^h})$.
\end{proof}

Formula \eqref{eq:L-layer} is a genuine $q$-series bridge between the
Pascal zero divisor and Lambert-series theory.  It is distinct from the
Dirichlet spectral zeta function \eqref{eq:spectral-zeta}.

\subsection{Rank recursion and a homogeneous Mahler equation}

Let $T$ be the Mahler operator $(Tf)(q)=f(q^2)$ and put
$R(q)=q/(1-q)$.

\begin{proposition}[Pascal--Mahler recursion]\label{prop:L-recursion}
The Lambert hierarchy satisfies
\begin{align}
 (I-T)L_1&=R,\label{eq:L-r1}\\
 (I-T)L_r&=TL_{r-1},\qquad r\geq2,\label{eq:L-rec}\\
 (I-T)^rL_r&=T^{r-1}R
 =\frac{q^{2^{r-1}}}{1-q^{2^{r-1}}}.\label{eq:L-iterated}
\end{align}
\end{proposition}

\begin{proof}
Subtract \eqref{eq:L-layer} evaluated at $q^2$ from the same formula at $q$.
Pascal's identity
$\binom h{r-1}-\binom{h-1}{r-1}=\binom{h-1}{r-2}$ gives
\eqref{eq:L-rec}; the rank-one difference is the first Lambert atom $R$.
Since $T$ commutes with $I-T$, induction gives \eqref{eq:L-iterated}.
\end{proof}

Let $a=2^{r-1}$.  The rational function
$g(q)=q^a/(1-q^a)$ satisfies
\[
 (1+q^a)g(q^2)-q^ag(q)=0.
\]
Combining this with \eqref{eq:L-iterated} gives the promised homogeneous
equation.

\begin{theorem}[Order-$(r+1)$ Mahler equation]\label{thm:L-Mahler}
For every $r\geq1$,
\begin{equation}\label{eq:L-homogeneous}
 \boxed{
 \bigl((1+q^{2^{r-1}})T-q^{2^{r-1}}I\bigr)
 (I-T)^rL_r=0.}
\end{equation}
In expanded form this is a nontrivial linear Mahler equation involving
$L_r(q),L_r(q^2),\ldots,L_r(q^{2^{r+1}})$ with rational coefficients.
\end{theorem}

\subsection{Dense dyadic singularities}

\begin{theorem}[Natural boundary of $L_r$]\label{thm:L-boundary}
Every dyadic root of unity is a radial singularity of $L_r$.  Since these
roots are dense on the unit circle, $|q|=1$ is a natural boundary.
\end{theorem}

\begin{proof}
Let $\zeta$ have order $2^J$ and put $q=\rho\zeta$, $0<\rho<1$.  For every
$h\geq\max\{J,r-1\}$,
\[
 q^{2^h}=\rho^{2^h}>0.
\]
Thus the tail of \eqref{eq:L-layer} is positive real.  Already the fixed term
with $h=\max\{J,r-1\}$ tends to $+\infty$ as $\rho\uparrow1$, while the
finitely many preceding terms cannot cancel that divergence.  Hence $\zeta$
is a radial singularity.  An analytic continuation through any open arc of
the circle would continue across some dyadic root, a contradiction.
\end{proof}

\begin{theorem}[Hypertranscendence and nonholonomicity of $L_r$]\label{thm:L-hyper}
For every $r\geq1$:
\begin{enumerate}[label=(\roman*)]
\item $L_r$ is hypertranscendental over $\ComplexNumbers(q)$;
\item $L_r$ is not $D$-finite;
\item the coefficient sequence $\mu_r(n)$ is not $P$-recursive.
\end{enumerate}
\end{theorem}

\begin{proof}
The natural boundary rules out rationality, including rationality in a
ramified variable.  Apply \cref{thm:ADH} to \eqref{eq:L-homogeneous}.
Hypertranscendence implies non-$D$-finiteness.  Equivalently, a $D$-finite
ordinary generating function has a $P$-recursive coefficient sequence, so
$\mu_r(n)$ is not $P$-recursive.
\end{proof}

\begin{remark}[Regular but not holonomic]
For fixed $r$, $\mu_r(n)$ is nevertheless $2$-regular.  With the auxiliary
convention $\mu_0(n):=1$,
\begin{equation}\label{eq:mu-2regular}
 \mu_r(2n)=\mu_r(n)+\mu_{r-1}(n),\qquad
 \mu_r(2n+1)=\begin{cases}1,&r=1,\\0,&r\geq2.\end{cases}
\end{equation}
Thus finite-state linear algebra controls its $2$-kernel even though no
polynomial-coefficient recurrence in $n$ exists.
\end{remark}

\subsection{The master rank Lambert germ}

\begin{theorem}[Bivariate rank resolvent]\label{thm:master-L}
For $|q|<1$ and $u$ in any bounded set for which the displayed series
converges normally,
\begin{equation}\label{eq:master-L}
 \mathscr L(u,q):=\sum_{r\geq1}u^{r-1}L_r(q)
 =\sum_{h\geq0}(1+u)^h\frac{q^{2^h}}{1-q^{2^h}}.
\end{equation}
It satisfies the first-order inhomogeneous Mahler equation
\begin{equation}\label{eq:master-L-Mahler}
 \boxed{\mathscr L(u,q)=\frac{q}{1-q}+(1+u)\mathscr L(u,q^2).}
\end{equation}
\end{theorem}

\begin{proof}
Interchange the rank and scale sums in \eqref{eq:L-layer} and use
\[
 \sum_{r=1}^{h+1}\binom h{r-1}u^{r-1}=(1+u)^h.
\]
Separating the $h=0$ term and shifting $h\mapsto h+1$ gives
\eqref{eq:master-L-Mahler}.
\end{proof}

\begin{remark}[A $q$-analogue viewpoint]
The factors $q^{2^h}/(1-q^{2^h})$ are logarithmic derivatives of dyadic
$q$-Pochhammer atoms.  The rank parameter changes the scale weight from $1$
to $(1+u)^h$.  This makes \eqref{eq:master-L} a natural meeting point of
automatic sequences, Lambert series, and the repository's broader
$q$-Pochhammer program.
\end{remark}

\section{Sharp higher-rank dyadic-comb cubature}

\subsection{A general periodization identity}

Let $Q\geq1$ be an integer lattice parameter; it is denoted $Q$ rather than
$q$ to avoid conflict with the Lambert variable.  For a polynomial $p$, put
\begin{equation}\label{eq:comb}
 \mathcal C_{r,Q}[p](x):=\frac{1}{Q}\sum_{k\in\IntegerNumbers}
 p\!\left(x-\frac{k}{Q}\right)
 f_r\!\left(x-\frac{k}{Q}\right).
\end{equation}
The sum is finite because $f_r$ is supported on $[-1,1]$.

\begin{lemma}[Fourier coefficients of the weighted comb]\label{lem:comb-Fourier}
For $p(t)=t^j$,
\begin{equation}\label{eq:comb-Fourier}
 \mathcal C_{r,Q}[t^j](x)
 =\sum_{m\in\IntegerNumbers}
 \frac{\Phi_r^{(j)}(Qm)}{(-2\pi\ImaginaryUnit)^j}
 \EulerE^{2\pi\ImaginaryUnit Qmx}.
\end{equation}
The constant coefficient is $m_{r,j}$.
\end{lemma}

\begin{proof}
Apply Poisson summation to $g_j(t)=t^jf_r(t)$ on the lattice $Q^{-1}\IntegerNumbers$.
Our Fourier convention gives
$\widehat g_j(\xi)=\Phi_r^{(j)}(\xi)/(-2\pi\ImaginaryUnit)^j$.
At $\xi=0$ this is $\int t^jf_r(t)\Differential t=m_{r,j}$.
\end{proof}

\subsection{Exact order and sharpness}

\begin{theorem}[Sharp Pascal comb theorem]\label{thm:sharp-comb}
Let
\begin{equation}\label{eq:M-rQ}
 M_r(Q):=\mu_r(Q)=\binom{\TwoAdicValuation(Q)+1}{r}.
\end{equation}
If $p$ is a polynomial with $\deg p<M_r(Q)$, then for every real $x$,
\begin{equation}\label{eq:comb-exact}
 \boxed{
 \frac1Q\sum_{k\in\IntegerNumbers}p\!\left(x-\frac{k}{Q}\right)
 f_r\!\left(x-\frac{k}{Q}\right)=\Expectation p(X_r).}
\end{equation}
The threshold is sharp.  If $M_r(Q)>0$, the identity fails for
$p(t)=t^{M_r(Q)}$; if $M_r(Q)=0$, it already fails for constants.
\end{theorem}

\begin{proof}
For $m\neq0$,
\[
 \mu_r(Qm)=\binom{\TwoAdicValuation(Q)+\TwoAdicValuation(m)+1}{r}\geq M_r(Q).
\]
Hence every derivative $\Phi_r^{(j)}(Qm)$ vanishes for $j<M_r(Q)$.  By
linearity and \cref{lem:comb-Fourier}, all nonconstant Fourier modes disappear
for every polynomial of degree below the threshold, leaving the moment at
frequency zero.

If $M=M_r(Q)>0$, the zero at $Q$ has exact order $M$, so
$\Phi_r^{(M)}(Q)\neq0$.  The $m=1$ Fourier coefficient of the periodization for
$t^M$ is therefore nonzero.  If $M=0$, then $\Phi_r(Q)\neq0$, so even the
constant periodization has a nonzero first Fourier mode.
\end{proof}

\begin{corollary}[Dyadic degree]\label{cor:dyadic-degree}
For $Q=2^H$, the exact polynomial degree is
\begin{equation}\label{eq:dyadic-degree}
 d_{r,H}=\binom{H+1}{r}-1
\end{equation}
when $H\geq r-1$; when $H<r-1$, even the constant periodization fails.
\end{corollary}

For example:
\begin{center}
\begin{tabular}{@{}c|rrrrrr@{}}
\toprule
& $H=2$ & $H=3$ & $H=4$ & $H=5$ & $H=6$ & $H=7$\\
\midrule
$r=1$ & 2 & 3 & 4 & 5 & 6 & 7\\
$r=2$ & 2 & 5 & 9 & 14 & 20 & 27\\
$r=3$ & 0 & 3 & 9 & 19 & 34 & 55\\
$r=4$ & -- & 0 & 4 & 14 & 34 & 69\\
\bottomrule
\end{tabular}
\end{center}
Here ``--'' means $M_r(2^H)=0$.

\begin{remark}[Relation to the known rank-one theorem]
At $r=1$, $M_1(Q)=1+\TwoAdicValuation(Q)$, exactly the sharp rank-one hierarchy already
proved in the dyadic-comb and frontier-compilation volumes.  The new content is
the Pascal-rank lift and its binomial growth in $\TwoAdicValuation(Q)$.
\end{remark}

\subsection{Appell reconstruction}

Let
\begin{equation}\label{eq:mgf}
 \Mr_r(z):=\Expectation\EulerE^{zX_r}=\Phi_r\!\left(\frac{\ImaginaryUnit z}{2\pi}\right).
\end{equation}
It is even.  Define Appell polynomials $P_{r,n}$ by
\begin{equation}\label{eq:Appell}
 \frac{\EulerE^{yz}}{\Mr_r(z)}
 =\sum_{n\geq0}P_{r,n}(y)\frac{z^n}{n!}.
\end{equation}

\begin{theorem}[Exact polynomial reproduction]\label{thm:Appell-reproduction}
For $0\leq n<M_r(Q)$,
\begin{equation}\label{eq:Appell-reproduce}
 \boxed{
 \frac1Q\sum_{k\in\IntegerNumbers}P_{r,n}\!\left(\frac{k}{Q}\right)
 f_r\!\left(x-\frac{k}{Q}\right)=x^n.}
\end{equation}
\end{theorem}

\begin{proof}
Apply \cref{thm:sharp-comb} to the polynomial in the variable $t$
\[
 p_x(t)=P_{r,n}(x-t).
\]
Its degree is $n<M_r(Q)$, so the left side of \eqref{eq:Appell-reproduce}
is $\Expectation P_{r,n}(x-X_r)$.  Taking the exponential generating function and using
symmetry,
\[
 \sum_{n\geq0}\Expectation P_{r,n}(x-X_r)\frac{z^n}{n!}
 =\frac{\EulerE^{xz}\Expectation\EulerE^{-zX_r}}{\Mr_r(z)}
 =\EulerE^{xz}.
\]
Coefficient comparison gives $x^n$.
\end{proof}

\subsection{Legendre and Bernoulli corollaries}

Let $\mathrm P_n$ denote the Legendre polynomial and $B_n(x)$ the Bernoulli
polynomial.

\begin{corollary}[Exact finite Legendre moments]\label{cor:Legendre-comb}
If $n<M_r(Q)$, then for every $x$,
\begin{equation}\label{eq:Legendre-comb}
 \int_{-1}^{1}\mathrm P_n(t)f_r(t)\Differential t
 =\frac1Q\sum_{|x-k/Q|\leq1}
 \mathrm P_n\!\left(x-\frac{k}{Q}\right)
 f_r\!\left(x-\frac{k}{Q}\right).
\end{equation}
\end{corollary}

\begin{corollary}[Exact finite Bernoulli moments]\label{cor:Bernoulli-comb}
If $n<M_r(Q)$, then
\begin{equation}\label{eq:Bernoulli-comb}
 \Expectation B_n(X_r)
 =\frac1Q\sum_{|x-k/Q|\leq1}
 B_n\!\left(x-\frac{k}{Q}\right)
 f_r\!\left(x-\frac{k}{Q}\right).
\end{equation}
\end{corollary}

These identities turn initial blocks of orthogonal and Appell moments into
finite exact sampling formulas.  They are potentially useful for stable
Legendre coefficient computation because the right side is local and the
exactness degree is known before any numerical work begins.

\subsection{Numerical zero-jet check}

The archived program differentiates truncated products at three exact zeros.
The first nonzero derivative agrees with
$M_r(Q)=\binom{\TwoAdicValuation(Q)+1}{r}$:
\begin{center}
\begin{tabular}{@{}rrrr@{}}
\toprule
$r$ & $Q$ & predicted order & first nonzero derivative (numerical)\\
\midrule
2 & 4  & 3 & $\Phi_2^{(3)}(4)\approx 1.32282664349\times10^{-2}$\\
3 & 8  & 4 & $\Phi_3^{(4)}(8)\approx-7.57646125312\times10^{-5}$\\
4 & 16 & 5 & $\Phi_4^{(5)}(16)\approx 3.51415041684\times10^{-8}$\\
\bottomrule
\end{tabular}
\end{center}
The lower derivatives are numerically below $10^{-130}$ at the chosen
precision.  This table checks the implementation; sharpness in
\cref{thm:sharp-comb} follows analytically from the exact zero order.

\section{All-rank Laguerre--P\'olya and P\'olya-frequency structure}

\subsection{Wick rotation}

From \eqref{eq:canonical},
\begin{equation}\label{eq:A-product}
 \Ar_r(w):=\Phi_r(\ImaginaryUnit\sqrt w)
 =\prod_{n\geq1}\left(1+\frac{w}{n^2}\right)^{\mu_r(n)}
 =\sum_{k\geq0}C_{r,k}w^k.
\end{equation}
Since $X_r$ is symmetric,
\begin{equation}\label{eq:A-mgf}
 \Ar_r(w)=\Expectation\EulerE^{2\pi\sqrt w X_r}=\Expectation\EulerE^{-2\pi\sqrt w X_r}
 =\sum_{k\geq0}\frac{(2\pi)^{2k}m_{r,2k}}{(2k)!}w^k.
\end{equation}
The square root disappears from the Taylor series.

\subsection{Classification and total positivity}

\begin{theorem}[All-rank type-I Laguerre--P\'olya theorem]\label{thm:LP-all}
For every $r\geq1$, $\Ar_r$ is a type-I Laguerre--P\'olya entire function of
genus zero and order $1/2$.  Its zeros are exactly $-n^2$, with multiplicity
$\mu_r(n)$.
\end{theorem}

\begin{proof}
The convergence
$\sum_n\mu_r(n)/n^2=\ZetaR_r(2)<\infty$ makes
\eqref{eq:A-product} a genus-zero product and the locally uniform limit of
polynomials with only nonpositive real zeros.  Hence it is type-I
Laguerre--P\'olya.  The exponent of convergence of the zeros is $1/2$ because
\[
 \sum_{n\geq1}\frac{\mu_r(n)}{(n^2)^\rho}=\ZetaR_r(2\rho)
\]
converges exactly for $\rho>1/2$.  Thus the order is at least $1/2$.
On the other hand, $|X_r|\leq1$ and \eqref{eq:A-mgf} give growth
$\exp(\BigO(\sqrt{|w|}))$, so the order is at most $1/2$.
\end{proof}

\begin{theorem}[$PF_\infty$ coefficients and Jensen hyperbolicity]\label{thm:PF-all}
For every $r\geq1$, the ordinary coefficient sequence
$(C_{r,k})_{k\geq0}$ is a P\'olya-frequency sequence of infinite order.
Equivalently, every minor of the Toeplitz matrix
$(C_{r,j-i})_{i,j\geq0}$ is nonnegative.  All shifted Jensen polynomials of
the exponential coefficient sequence $k!C_{r,k}$ are hyperbolic with
nonpositive zeros.
\end{theorem}

\begin{proof}
The Aissen--Edrei--Schoenberg--Whitney characterization of $PF_\infty$
sequences applies to
\[
 \Ar_r(w)=\prod_{n,j}\left(1+\frac{w}{n^2}\right),
 \qquad 1\leq j\leq\mu_r(n),
\]
because all parameters $1/n^2$ are nonnegative and summable
\cite{aesw}.  Jensen's characterization of the Laguerre--P\'olya class then
gives the hyperbolicity assertion.
\end{proof}

\begin{corollary}[Ultra-log-concavity of every rank's moments]\label{cor:ULC}
For $k\geq1$,
\begin{align}
 C_{r,k}^2&\geq\frac{k+1}{k}C_{r,k-1}C_{r,k+1},\label{eq:C-ULC}\\
 m_{r,2k}^2&\geq\frac{2k-1}{2k+1}
 m_{r,2k-2}m_{r,2k+2}.\label{eq:moment-ULC}
\end{align}
\end{corollary}

\begin{proof}
Apply Newton's inequalities to finite truncations of
\eqref{eq:A-product} and pass to the locally uniform limit.  Substitution of
\eqref{eq:C-moment} simplifies the factorial ratio to
$(2k-1)/(2k+1)$.
\end{proof}

The numerical ratios
\[
 \frac{C_{r,k}^2}{((k+1)/k)C_{r,k-1}C_{r,k+1}}
\]
are all greater than $1$ in the archived table for $1\leq r\leq6$ and
$1\leq k\leq6$.  For example, the $k=1$ ratios are approximately
$1.31579,1.16822,1.09457,1.05467,1.03210$ for ranks $1$ through $5$.

\begin{remark}[Repository boundary]
The rank-one instance, including the $PF_\infty$ theorem and
\eqref{eq:moment-ULC}, already appears in
\path{Frontier_Compilations.tex}.  The theorem here is the all-rank lift through
the Pascal multiplicities.
\end{remark}

\subsection{Rank convolution}

For $r\geq2$, Wick rotation of \eqref{eq:refinement} gives
\begin{equation}\label{eq:A-refinement}
 \Ar_r(4w)=\Ar_r(w)\Ar_{r-1}(w).
\end{equation}

\begin{corollary}[Coefficient Pascal convolution]\label{cor:C-convolution}
For $r\geq2$ and every $k\geq0$,
\begin{equation}\label{eq:C-convolution}
 4^kC_{r,k}=\sum_{j=0}^{k}C_{r,j}C_{r-1,k-j}.
\end{equation}
\end{corollary}

This identity may be read either as a moment convolution or as a positive
coefficient recursion inside the $PF_\infty$ cone.

\subsection{Non-\texorpdfstring{$D$}{D}-finiteness of the Fourier products}

\begin{theorem}[Unbounded zero orders obstruct holonomicity]\label{thm:Phi-nonD}
For every $r\geq1$, neither $\Phi_r(z)$ nor $\Ar_r(w)$ is $D$-finite over
$\ComplexNumbers(z)$.
\end{theorem}

\begin{proof}
Suppose a nonzero entire function $y$ satisfies a linear differential equation
of order $d$ with polynomial coefficients.  Away from the finitely many zeros
of the leading coefficient, an ordinary-point uniqueness theorem says that a
zero of order at least $d$ forces the solution to vanish identically.  Thus
zero multiplicities at ordinary points are uniformly bounded by $d-1$.

For $\Phi_r$, the points $2^H$ have multiplicity
$\binom{H+1}{r}\to\infty$; for $\Ar_r$, the points $-4^H$ have the same
unbounded multiplicities.  Infinitely many of those points avoid the finite
singular set, a contradiction.
\end{proof}

\begin{conjecture}[Differential transcendence of the sinc hierarchy]\label{conj:Phi-hyper}
For every $r\geq1$, $\Phi_r$ and $\Ar_r$ are hypertranscendental over
$\ComplexNumbers(z)$.  More strongly, finite sets of distinct ranks should be
differentially algebraically independent after quotienting by the explicit
refinement relations.
\end{conjecture}

The Mahler dichotomy used for $G_r$ and $L_r$ does not apply directly here:
\eqref{eq:refinement} couples different ranks and is multiplicative rather
than a scalar linear Mahler equation with rational coefficients.

\section{The Poisson-binomial spectral count}

\subsection{Exact probability generating function}

Fix $r\geq1$ and $\theta>0$.  For every pair $(n,j)$ with
$1\leq j\leq\mu_r(n)$, let $B_{n,j}$ be independent Bernoulli variables with
\begin{equation}\label{eq:p-n}
 p_{n,\theta}:=\Probability(B_{n,j}=1)=\frac{\theta}{n^2+\theta}.
\end{equation}
Since
$\sum_{n,j}p_{n,\theta}\leq\theta\ZetaR_r(2)<\infty$, the sum below is finite
almost surely.

\begin{definition}[Spectral count]\label{def:N}
\begin{equation}\label{eq:N-def}
 N_{r,\theta}:=\sum_{n\geq1}\sum_{j=1}^{\mu_r(n)}B_{n,j}.
\end{equation}
\end{definition}

\begin{theorem}[Poisson-binomial determinant duality]\label{thm:PB-dual}
The probability generating function of $N_{r,\theta}$ is
\begin{equation}\label{eq:PB-pgf}
 \boxed{
 \Expectation z^{N_{r,\theta}}=
 \frac{\Ar_r(\theta z)}{\Ar_r(\theta)}.}
\end{equation}
Consequently,
\begin{equation}\label{eq:PB-prob}
 \Probability(N_{r,\theta}=k)
 =\frac{C_{r,k}\theta^k}{\Ar_r(\theta)}
 =\frac{(2\pi)^{2k}m_{r,2k}\theta^k}
 {(2k)!\Ar_r(\theta)}.
\end{equation}
\end{theorem}

\begin{proof}
Each factor in \eqref{eq:A-product} normalizes as
\[
 \frac{1+\theta z/n^2}{1+\theta/n^2}
 =\frac{n^2}{n^2+\theta}+\frac{\theta}{n^2+\theta}z,
\]
the PGF of a Bernoulli variable with parameter \eqref{eq:p-n}.  Multiply the
factors with their multiplicities.  Coefficient extraction and
\eqref{eq:C-moment} give \eqref{eq:PB-prob}.
\end{proof}

At $\theta=1$, the normalized even-moment sequence is literally a probability
mass function:
\begin{equation}\label{eq:moment-pmf}
 \Probability(N_{r,1}=k)=
 \frac{(2\pi)^{2k}m_{r,2k}}{(2k)!\Ar_r(1)}.
\end{equation}
This is an exact moment--spectrum probability dual, not merely an analogy.

\subsection{Cumulants, branching, and coefficient tilting}

\begin{proposition}[Mean, variance, and factorial cumulants]\label{prop:PB-cumulants}
Let
\begin{equation}\label{eq:S-k}
 S_{r,k}(\theta):=\sum_{n\geq1}\mu_r(n)
 \left(\frac{\theta}{n^2+\theta}\right)^k.
\end{equation}
Then
\begin{align}
 \lambda_{r,\theta}:=\Expectation N_{r,\theta}&=S_{r,1}(\theta),\label{eq:PB-mean}\\
 \Variance(N_{r,\theta})&=S_{r,1}(\theta)-S_{r,2}(\theta),\label{eq:PB-var}\\
 \kappa_k^{(F)}(N_{r,\theta})
 &=(-1)^{k-1}(k-1)!S_{r,k}(\theta).\label{eq:PB-factorial}
\end{align}
\end{proposition}

\begin{proof}
Use $\log\Expectation(1+t)^N=\sum_{n,j}\log(1+p_{n,\theta}t)$ and compare powers of
$t$.
\end{proof}

\begin{theorem}[Rank branching]\label{thm:N-branching}
For $r\geq2$,
\begin{equation}\label{eq:N-branching}
 N_{r,4\theta}\stackrel d=N_{r,\theta}+N_{r-1,\theta}',
\end{equation}
with independent variables on the right.
\end{theorem}

\begin{proof}
Normalize \eqref{eq:A-refinement} at $w=\theta$ and at $w=\theta z$; the
resulting PGF factorization is exactly convolution of independent counts.
\end{proof}

Equation \eqref{eq:PB-prob} also gives an exact exponential tilt for
coefficient extraction.

\begin{corollary}[Saddle tilt]\label{cor:saddle-tilt}
For every $\theta>0$,
\begin{equation}\label{eq:saddle-tilt}
 C_{r,k}=\frac{\Ar_r(\theta)}{\theta^k}
 \Probability(N_{r,\theta}=k).
\end{equation}
The map $\theta\mapsto\lambda_{r,\theta}$ is strictly increasing from $0$ to
$\infty$, so every $k>0$ has a unique saddle parameter satisfying
$\lambda_{r,\theta}=k$.
\end{corollary}

\begin{proof}
The first identity is a rearrangement of \eqref{eq:PB-prob}.  Each Bernoulli
parameter is strictly increasing in $\theta$, so the mean is strictly
increasing.  Monotone convergence gives the endpoint limits; as
$\theta\to\infty$, increasingly many spectral levels have probability bounded
away from zero.
\end{proof}

This converts asymptotics of the moment coefficients into local limit theory
for $N_{r,\theta}$, and conversely.

\subsection{Basic comparison inequalities}

The following estimates drive the rank phase diagram.

\begin{lemma}[Zeta comparison]\label{lem:PB-bounds}
For every $r$ and $\theta>0$,
\begin{align}
 0&\leq \theta\ZetaR_r(2)-\lambda_{r,\theta}
 \leq\theta^2\ZetaR_r(4),\label{eq:mean-bound}\\
 S_{r,2}(\theta)&\leq\theta^2\ZetaR_r(4).\label{eq:S2-bound}
\end{align}
\end{lemma}

\begin{proof}
Termwise,
\[
 0\leq\frac{\theta}{n^2}-\frac{\theta}{n^2+\theta}
 =\frac{\theta^2}{n^2(n^2+\theta)}\leq\frac{\theta^2}{n^4},
\]
and
$(\theta/(n^2+\theta))^2\leq\theta^2/n^4$.
Sum with multiplicities.
\end{proof}

\subsection{Critical Poisson limit}

\begin{theorem}[Critical rank Poisson law]\label{thm:critical-Poisson}
Fix $c>0$ and put $\theta_r=c3^r$.  Then
\begin{equation}\label{eq:critical-Poisson}
 N_{r,c3^r}\xrightarrow[r\to\infty]{\mathrm{TV}}
 \Po\!\left(4c\zeta(2)\right)
 =\Po\!\left(\frac{2c\pi^2}{3}\right).
\end{equation}
With $\TotalVariationDistance$ defined as the supremum over events,
\begin{equation}\label{eq:critical-TV}
 \TotalVariationDistance\!\left(N_{r,c3^r},\Po(4c\zeta(2))\right)
 \leq48c^2\zeta(4)\left(\frac35\right)^r.
\end{equation}
\end{theorem}

\begin{proof}
By \eqref{eq:zeta-special},
$\theta_r\ZetaR_r(2)=4c\zeta(2)$.  The mean error is at most
\[
 \theta_r^2\ZetaR_r(4)
 =16c^2\zeta(4)(3/5)^r.
\]
Le Cam's inequality for a Poisson-binomial sum gives
\[
 \TotalVariationDistance(N_{r,\theta_r},\Po(\lambda_{r,\theta_r}))
 \leq2S_{r,2}(\theta_r)
 \leq2\theta_r^2\ZetaR_r(4).
\]
The total variation distance between Poisson laws with parameters $\lambda$
and $\lambda'$ is at most $|\lambda-\lambda'|$.  Add the two bounds.
\end{proof}

For $c=1$, the limiting mean is
$4\zeta(2)=6.5797362673929057\ldots$.  The grouped numerical means are:
\begin{center}
\begin{tabular}{@{}rrrr@{}}
\toprule
$r$ & $\Expectation N_{r,3^r}$ & target minus mean & proved TV bound\\
\midrule
8  & 6.343771060790705 & $2.35965\times10^{-1}$ & $8.72586\times10^{-1}$\\
12 & 6.544298628120802 & $3.54376\times10^{-2}$ & $1.13087\times10^{-1}$\\
16 & 6.574932579970129 & $4.80369\times10^{-3}$ & $1.46561\times10^{-2}$\\
20 & 6.579105941567851 & $6.30326\times10^{-4}$ & $1.89943\times10^{-3}$\\
24 & 6.579654307729619 & $8.19597\times10^{-5}$ & $2.46166\times10^{-4}$\\
28 & 6.579725636253052 & $1.06311\times10^{-5}$ & $3.19031\times10^{-5}$\\
32 & 6.579734889287106 & $1.37811\times10^{-6}$ & $4.13465\times10^{-6}$\\
\bottomrule
\end{tabular}
\end{center}

\subsection{A Poisson, mod-Poisson, and Gaussian phase diagram}

The isolated critical result extends to a multi-scale theorem.

\begin{theorem}[Rank phase diagram]\label{thm:PB-phase}
Fix $c>0$, $a>0$, and set $\theta_r=ca^r$.  Let
$\lambda_r=\Expectation N_{r,\theta_r}$ and
$\sigma_r^2=\Variance(N_{r,\theta_r})$.

\begin{enumerate}[label=(\roman*)]
\item If $0<a<3$, then $N_{r,\theta_r}\to0$ in probability and in total
variation.

\item If $a=3$, the fixed Poisson limit of
\cref{thm:critical-Poisson} holds.

\item If $3<a<\sqrt{15}$, then
\begin{equation}\label{eq:diverging-Poisson}
 \TotalVariationDistance(N_{r,\theta_r},\Po(\lambda_r))\longrightarrow0,
 \qquad
 \lambda_r\sim4c\zeta(2)(a/3)^r\longrightarrow\infty.
\end{equation}

\item If $a=\sqrt{15}$, then $\lambda_r\to\infty$ and the following
mod-Poisson convergence holds locally uniformly for $t\in\ComplexNumbers$:
\begin{equation}\label{eq:mod-Poisson}
 \Expectation(1+t)^{N_{r,\theta_r}}\EulerE^{-\lambda_rt}
 \longrightarrow
 \exp\!\left(-8c^2\zeta(4)t^2\right).
\end{equation}

\item Throughout the larger range $3<a<5$,
\begin{equation}\label{eq:PB-CLT}
 \frac{N_{r,\theta_r}-\lambda_r}{\sigma_r}
 \xrightarrow{d}\Normal(0,1),
\end{equation}
with
\begin{equation}\label{eq:mean-var-asymp}
 \lambda_r\sim\sigma_r^2\sim4c\zeta(2)(a/3)^r.
\end{equation}
Moreover a Berry--Esseen bound is
\begin{equation}\label{eq:PB-BE}
 \sup_x\left|\Probability\!\left(
 \frac{N_{r,\theta_r}-\lambda_r}{\sigma_r}\leq x\right)-\Phi_{\rm G}(x)
 \right|=\BigO((3/a)^{r/2}),
\end{equation}
where $\Phi_{\rm G}$ is the standard normal CDF.
\end{enumerate}
\end{theorem}

\begin{proof}
If $a<3$, Markov's inequality and
$\lambda_r\leq\theta_r\ZetaR_r(2)=4c\zeta(2)(a/3)^r$ prove (i).
Part (ii) was proved above.

For (iii), Le Cam and \eqref{eq:S2-bound} give
\[
 \TotalVariationDistance(N_{r,\theta_r},\Po(\lambda_r))
 \leq32c^2\zeta(4)(a^2/15)^r\to0.
\]
The mean asymptotic follows from
\eqref{eq:mean-bound}, because the relative error is bounded by a constant
times $(a/5)^r$.

Now let $a=\sqrt{15}$.  The maximum Bernoulli parameter tends to zero because
the smallest spectral point is $2^{r-1}$ and
\[
 \max_np_{n,\theta_r}\leq4c(a/4)^r\to0.
\]
Furthermore
\[
 0\leq\theta_r^2\ZetaR_r(4)-S_{r,2}(\theta_r)
 \leq2\theta_r^3\ZetaR_r(6)\to0,
\]
since $15^{3/2}/63<1$.  Thus
$S_{r,2}(\theta_r)\to16c^2\zeta(4)$.  For every $k\geq3$,
\[
 S_{r,k}(\theta_r)\leq\theta_r^k\ZetaR_r(2k)
 =2^{2k}c^k\zeta(2k)
 \left(\frac{15^{k/2}}{4^k-1}\right)^r\to0.
\]
Expanding
$\log\Expectation(1+t)^N=\sum_{k\geq1}(-1)^{k-1}S_{r,k}t^k/k$ proves
\eqref{eq:mod-Poisson}; the vanishing maximum parameter makes the remainder
uniform on compact $t$-sets.

Finally suppose $3<a<5$.  By \cref{lem:PB-bounds},
\[
 \lambda_r=\theta_r\ZetaR_r(2)(1+\LittleO(1)),\qquad
 S_{r,2}(\theta_r)=\LittleO(\lambda_r),
\]
so \eqref{eq:mean-var-asymp} holds and $\sigma_r\to\infty$.  For a Bernoulli
variable $B$ with parameter $p$,
$\Expectation|B-p|^3\leq p(1-p)$.  Hence the sum of third absolute centered moments is
at most $\sigma_r^2$.  The Berry--Esseen theorem, first for finite truncations
and then by an $L^2$ tail limit, gives an error $\BigO(1/\sigma_r)$, which is
\eqref{eq:PB-BE}.
\end{proof}

\begin{remark}[Two distinct critical scales]
The scale $3^r$ comes from the first spectral trace $\ZetaR_r(2)$ and creates
the fixed Poisson law.  The scale $(\sqrt{15})^r$ makes the second factorial
cumulant order one and creates the mod-Poisson correction.  The relative
variance correction becomes order one only at the still larger scale $5^r$.
This hierarchy of critical bases is a direct probabilistic shadow of the
spectral denominators $(2^{2k}-1)^r$.
\end{remark}

\section{A generalized-gamma spectral dual and its rank CLT}

\subsection{The Thorin measure}

Fix $\tau>0$.  Let
\begin{equation}\label{eq:V-def}
 V_{r,\tau}:=\sum_{n\geq1}Y_{r,n},
 \qquad
 Y_{r,n}\sim\operatorname{Gamma}(\tau\mu_r(n),\text{ rate }n^2),
\end{equation}
with independent summands and zero-shape terms omitted.  Since
$\sum_n\Expectation Y_{r,n}=\tau\ZetaR_r(2)<\infty$, the sum is finite almost surely.

\begin{theorem}[Generalized-gamma determinant dual]\label{thm:gamma-dual}
For $w\geq0$,
\begin{equation}\label{eq:gamma-LT}
 \boxed{
 \Expectation\EulerE^{-wV_{r,\tau}}=\Ar_r(w)^{-\tau}.}
\end{equation}
Thus $V_{r,\tau}$ is a generalized gamma convolution with discrete Thorin
measure
\begin{equation}\label{eq:Thorin}
 \mathcal T_{r,\tau}=\tau\sum_{n\geq1}\mu_r(n)\delta_{n^2}.
\end{equation}
Its cumulants are
\begin{equation}\label{eq:gamma-cumulants}
 \kappa_k(V_{r,\tau})=\tau(k-1)!\ZetaR_r(2k),\qquad k\geq1.
\end{equation}
\end{theorem}

\begin{proof}
The Laplace transform of a gamma variable with shape $\alpha$ and rate
$\beta$ is $(1+w/\beta)^{-\alpha}$.  Multiply the transforms and use
\eqref{eq:A-product}.  Logarithmic differentiation gives
\eqref{eq:gamma-cumulants}.
\end{proof}

\begin{theorem}[Gamma rank branching]\label{thm:V-branching}
For $r\geq2$,
\begin{equation}\label{eq:V-branching}
 4V_{r,\tau}\stackrel d=V_{r,\tau}+V_{r-1,\tau}',
\end{equation}
with independence on the right.
\end{theorem}

\begin{proof}
Evaluate the Laplace transform of $4V_{r,\tau}$ at $w$ and apply
\eqref{eq:A-refinement}.
\end{proof}

\subsection{Concentration}

From \eqref{eq:zeta-special},
\begin{equation}\label{eq:V-mean-var}
 \Expectation V_{r,\tau}=\frac{4\tau\zeta(2)}{3^r},
 \qquad
 \Variance(V_{r,\tau})=\frac{16\tau\zeta(4)}{15^r}.
\end{equation}
Therefore
\begin{equation}\label{eq:V-relative}
 \frac{\Variance(V_{r,\tau})}{(\Expectation V_{r,\tau})^2}
 =\frac{2}{5\tau}\left(\frac35\right)^r.
\end{equation}

\begin{corollary}[Deterministic mean normalization]\label{cor:V-L2}
\begin{equation}\label{eq:V-L2}
 \frac{V_{r,\tau}}{\Expectation V_{r,\tau}}\longrightarrow1
 \qquad\text{in }L^2.
\end{equation}
\end{corollary}

The next result studies the much finer fluctuation scale.

\subsection{A rank central limit theorem}

Let
\begin{equation}\label{eq:W-def}
 W_{r,\tau}:=
 \frac{V_{r,\tau}-\Expectation V_{r,\tau}}
 {\sqrt{\Variance(V_{r,\tau})}}.
\end{equation}

\begin{theorem}[Generalized-gamma rank CLT]\label{thm:gamma-CLT}
For every fixed $\tau>0$,
\begin{equation}\label{eq:gamma-CLT}
 W_{r,\tau}\xrightarrow{d}\Normal(0,1).
\end{equation}
For $k\geq3$, its standardized cumulants are exactly
\begin{equation}\label{eq:gamma-standardized}
 \kappa_k(W_{r,\tau})
 =\tau^{1-k/2}(k-1)!
 \frac{\zeta(2k)}{\zeta(4)^{k/2}}
 \left(\frac{15^{k/2}}{4^k-1}\right)^r,
\end{equation}
and hence decay geometrically.
\end{theorem}

\begin{proof}
Formula \eqref{eq:gamma-standardized} follows directly from
\eqref{eq:gamma-cumulants} and \eqref{eq:V-mean-var}.  For every $k\geq3$,
$15^{k/2}<4^k-1$, so each fixed standardized cumulant tends to zero.

For a direct characteristic-function proof, put
$\sigma_r^2=\Variance(V_{r,\tau})$.  Then
\begin{align*}
 \log\Expectation\EulerE^{\ImaginaryUnit tW_{r,\tau}}
 &=-\frac{\ImaginaryUnit t\Expectation V_{r,\tau}}{\sigma_r}
 -\tau\sum_n\mu_r(n)
 \log\!\left(1-\frac{\ImaginaryUnit t}{\sigma_r n^2}\right).
\end{align*}
The smallest active $n$ is $2^{r-1}$, and
\begin{equation}\label{eq:gamma-radius}
 \sigma_r2^{2r-2}
 =\sqrt{\tau\zeta(4)}\left(\frac4{\sqrt{15}}\right)^r
 \longrightarrow\infty.
\end{equation}
Thus the logarithm may be expanded uniformly for fixed $t$.  The linear term
cancels the centering, and the quadratic term is $-t^2/2$.  The absolute
remainder is bounded by a constant multiple of
\[
 |t|^3\frac{\tau\ZetaR_r(6)}{\sigma_r^3}
 =\BigO\!\left(
 \left(\frac{15^{3/2}}{63}\right)^r\right)\to0.
\]
L\'evy's continuity theorem proves the CLT.
\end{proof}

The numerical standardized cumulants for $\tau=1$ are:
\begin{center}
\begin{tabular}{@{}rrrrr@{}}
\toprule
$r$ & $\kappa_3$ & $\kappa_4$ & $\kappa_5$ & $\kappa_6$\\
\midrule
8  & $9.4478\times10^{-1}$ & $1.8895$ & $5.4648$ & $2.0155\times10^{1}$\\
16 & $4.9397\times10^{-1}$ & $6.9420\times10^{-1}$ & $1.5149$ & $4.2908$\\
24 & $2.5827\times10^{-1}$ & $2.5505\times10^{-1}$ & $4.1996\times10^{-1}$ & $9.1347\times10^{-1}$\\
32 & $1.3503\times10^{-1}$ & $9.3704\times10^{-2}$ & $1.1642\times10^{-1}$ & $1.9447\times10^{-1}$\\
48 & $3.6914\times10^{-2}$ & $1.2648\times10^{-2}$ & $8.9469\times10^{-3}$ & $8.8140\times10^{-3}$\\
64 & $1.0091\times10^{-2}$ & $1.7073\times10^{-3}$ & $6.8756\times10^{-4}$ & $3.9947\times10^{-4}$\\
\bottomrule
\end{tabular}
\end{center}

\section{Connections to Bell, Bernoulli, Legendre, \texorpdfstring{$q$}{q}-series, and inverse problems}

\subsection{Exact Bell--Bernoulli moment calculus}

Combining \eqref{eq:Bell-coeff} and \eqref{eq:C-moment} gives the explicit
closed formula
\begin{equation}\label{eq:moment-Bell}
 m_{r,2n}=\frac{(2n)!}{(2\pi)^{2n}n!}
 \Bell_n\!\left(
 \ZetaR_r(2),-1!\ZetaR_r(4),\ldots,
 (-1)^{n+1}(n-1)!\ZetaR_r(2n)
 \right).
\end{equation}
Substitution of Euler's zeta values converts this to a rational expression in
$B_2,B_4,\ldots,B_{2n}$ and the denominators
$(2^{2k}-1)^r$.  The first exact moments, also reproduced by the Python
program, include
\begin{align*}
 m_{1,2}&=\frac19,&m_{1,4}&=\frac{19}{675},&m_{1,6}&=\frac{583}{59535},\\
 m_{2,2}&=\frac1{27},&m_{2,4}&=\frac{107}{30375},&m_{2,6}&=\frac{27953}{56260575},\\
 m_{3,2}&=\frac1{81},&m_{3,4}&=\frac{571}{1366875},&m_{3,6}&=\frac{1165663}{53166243375}.
\end{align*}

The new $PF_\infty$ theorem adds a family of inequalities to this exact
arithmetic data.  Every Toeplitz minor becomes a polynomial inequality among
these Bell--Bernoulli expressions.  This is stronger than the first quadratic
Tur\'an inequality and gives a systematic target for symbolic verification.

\subsection{Legendre data from exact combs}

The repository already gives a Fourier--Legendre expansion of the classical
up-function.  The higher-rank comb theorem contributes a complementary local
route.  For $n<M_r(Q)$, the Legendre coefficient
\begin{equation}\label{eq:Legendre-coeff}
 \ell_{r,n}:=\frac{2n+1}{2}\int_{-1}^{1}f_r(t)\mathrm P_n(t)\Differential t
\end{equation}
can be computed from any phase $x$ by the finite sum in
\eqref{eq:Legendre-comb}.  Varying $x$ provides exact internal consistency
checks.  Choosing $Q=2^H$ makes the exactness degree grow like
$H^r/r!$, a dramatic higher-rank gain over the linear rank-one degree.

A promising computational program is to combine this exact low-degree block
with the repository's global Legendre tail estimates.  The comb handles the
initial coefficients without quadrature error; only the tail needs analytic
bounding.

\subsection{Mahler and \texorpdfstring{$q$}{q}-series bridges}

There are now two different Mahler mechanisms in the same hierarchy.

\begin{itemize}
\item The sign functions $G_r$ are products of binary cyclotomic-like factors
and satisfy a first-order homogeneous Mahler equation at the block base
$B_r=2^{P_r}$.
\item The multiplicity series $L_r$ are weighted dyadic Lambert series.  Their
rank recursion produces an order-$(r+1)$ homogeneous equation at base $2$.
\end{itemize}

Both are hypertranscendental, but for different structural reasons: $G_r$
needs a digit-character nonperiodicity argument, while $L_r$ has a direct
dense set of radial singularities.  The master germ \eqref{eq:master-L}
adds a rank variable to the $q$-series side and may be viewed as a resolvent of
the scale-shift operator.

\subsection{Inverse Fabius and Lambert-\texorpdfstring{$W$}{W} directions}

The primary repository exposition already develops endpoint and inverse
asymptotics with exact lower-Lambert coordinates and periodic corrections.
The results here suggest three nonduplicative continuations.

First, let $Q_r(u)$ be the quantile function of $X_r$.  Then
\[
 \Expectation p(X_r)=\int_0^1p(Q_r(u))\Differential u.
\]
Combining this with \eqref{eq:comb-exact} transports every polynomial below
the sharp threshold into an exact identity between a quantile integral and a
finite density comb.  For $r=1$, translating between $X_1$, the up-density,
and the bounded Fabius CDF turns this into exact integral constraints on the
inverse Fabius function.

Second, the saddle identity \eqref{eq:saddle-tilt} reduces coefficient
asymptotics for $C_{r,k}$ to inversion of
$\lambda_{r,\theta}=k$.  In joint large-$(r,k)$ regimes, the dyadic pole
lattice of the spectral zeta should produce logarithmically periodic
corrections.  Determining whether the nonlinear inversion is most naturally
expressed in a Lambert-$W$ coordinate is an open problem rather than a claim
of this report.

Third, the mod-Poisson scale $\theta_r=c(\sqrt{15})^r$ is governed by the
second spectral trace, whereas endpoint inversion is governed by Mellin poles
at the imaginary dyadic lattice.  A two-parameter saddle analysis could reveal
a common periodic phase connecting these apparently different Lambert
phenomena.

\section{Conjectures and frontier research program}

\subsection{Differential algebra}

\begin{conjecture}[Hypertranscendence of the Fourier hierarchy]
For each $r\geq1$, the entire functions $\Phi_r$ and $\Ar_r$ are
hypertranscendental over $\ComplexNumbers(z)$.
\end{conjecture}

\begin{conjecture}[Cross-rank differential independence]
After adjoining the explicit refinement relations
$\Phi_r(2z)=\Phi_r(z)\Phi_{r-1}(z)$, no further algebraic differential
relations exist among a finite set of ranks.
\end{conjecture}

A possible route is parametrized difference Galois theory for the triangular
rank system rather than for a single scalar Mahler equation.

\subsection{Strict total positivity and zero interlacing}

\begin{conjecture}[Strictness]
After removing the forced degeneracies caused by repeated spectral zeros,
all admissible Toeplitz minors of $(C_{r,k})$ are strictly positive, and the
zeros of consecutive nontrivial Jensen polynomials strictly interlace.
\end{conjecture}

Finite products reduce this to strict total positivity for elementary
symmetric functions with repeated parameters.  The difficulty is formulating
the exact admissibility condition that survives the infinite limit.

\subsection{Beyond the proved Poisson-binomial phase range}

Theorem \ref{thm:PB-phase} identifies the bases $3$, $\sqrt{15}$, and $5$.
Several further transitions should occur.

\begin{problem}[The $5^r$ variance-critical law]
Determine the limiting centered law when $\theta_r=c5^r$.  At this scale the
second Bernoulli power sum is of the same order as the mean, so the variance is
no longer asymptotic to the mean.  Higher factorial cumulants should decide
whether the correct limit is still Gaussian after a modified normalization.
\end{problem}

\begin{problem}[Approach to the local saturation scale $4^r$]
The largest Bernoulli parameter changes from vanishing to order one near
$\theta_r\asymp4^r$, even though the global CLT survives beyond that scale.
Separate the finitely many saturated low levels from the diffuse high-level
cloud and derive a convolution limit.
\end{problem}

\begin{conjecture}[Full mod-Poisson hierarchy]
For each $k\geq2$, the critical base
$a_k=(4^k-1)^{1/k}$ produces a mod-Poisson correction whose first surviving
term is the $k$th factorial cumulant.  The bases increase to $4$, creating an
accumulating hierarchy of local spectral transitions.
\end{conjecture}

The statement needs care because lower factorial cumulants may already
diverge; the correct renormalization should subtract all preceding terms.

\subsection{Gamma Edgeworth theory}

\begin{problem}[Explicit Edgeworth expansion]
Use \eqref{eq:gamma-standardized} to derive an all-orders Edgeworth expansion
for $W_{r,\tau}$ with a remainder uniform in a growing central window.  The
smallest geometric decay factor is
$15^{3/2}/63$, so the skewness term should control the first correction.
\end{problem}

The characteristic-function proof already gives an expanding analytic disk
of radius proportional to $(4/\sqrt{15})^r$, an unusually favorable setting
for uniform inversion.

\subsection{Exact arithmetic and formalization}

\begin{problem}[Denominator geometry]
Combine the Bell formula \eqref{eq:moment-Bell}, von Staudt--Clausen, and the
rank denominators $(2^{2k}-1)^r$ to determine exact prime valuations of
$m_{r,2n}$.  The $PF_\infty$ inequalities may impose additional divisibility
or sign constraints on reduced numerators.
\end{problem}

\begin{problem}[Machine-checked Mahler frontier]
Formalize the nonperiodicity theorem, the Lambert operator recurrence, and the
sharp comb theorem.  The analytic imports needed for full hypertranscendence
are substantial, but most of the new finite combinatorics and exact Fourier
zero-order reasoning can be separated from them.
\end{problem}

\section{Formalization roadmap}

A realistic Lean development can be staged so that each layer has independent
value.

\subsection{Tier I: finite combinatorics}

\begin{enumerate}
\item Define $\sigma_{r,j}$ and $\varepsilon_r(N)$ from binary digits.
\item Prove periodicity of the mask by Lucas's theorem.
\item Prove strong $B_r$-multiplicativity.
\item Formalize \cref{lem:eventual-pure} and the complement-bit identity
\eqref{eq:complement-sign}.
\item Conclude non-eventual-periodicity.
\item Prove the layer-cake identity \eqref{eq:mu-layer} and the coefficient
recurrences \eqref{eq:L-r1}--\eqref{eq:L-iterated} in formal power series.
\end{enumerate}

These steps require no complex analysis.

\subsection{Tier II: exact products and Fourier algebra}

\begin{enumerate}
\item Reuse the repository's weighted-scale multiplicity theorem to identify
zero orders.
\item Formalize the periodized Fourier coefficient calculation
\eqref{eq:comb-Fourier} under the existing Poisson-summation API.
\item Prove the exactness and sharpness theorem \ref{thm:sharp-comb}.
\item Define the Appell family and prove \eqref{eq:Appell-reproduce} by finite
polynomial algebra plus moment evaluation.
\item Formalize the product identity for $\Ar_r$ and coefficient convolution
\eqref{eq:C-convolution}.
\end{enumerate}

\subsection{Tier III: probability}

\begin{enumerate}
\item Construct countable independent Bernoulli variables with summable
parameters and identify the PGF \eqref{eq:PB-pgf}.
\item Prove the mean and power-sum bounds in \cref{lem:PB-bounds}.
\item Import or prove a countable version of Le Cam's inequality by finite
truncation.
\item Establish the critical Poisson theorem and the subcritical regime.
\item Formalize Berry--Esseen for finite truncations and pass to the limit.
\item Construct the gamma convolution from summable means and prove its
Laplace transform and cumulants.
\end{enumerate}

\subsection{Tier IV: external transcendence theorems}

The final hypertranscendence statements depend on the deep external theorem
\ref{thm:ADH}.  A practical formal boundary would encode that theorem as an
explicit imported assumption, prove that $G_r$ and $L_r$ satisfy its
hypotheses, and keep the elementary nonrationality and natural-boundary proofs
independent.

\section{Computational experiments and reproducibility}

\subsection{What the program checks}

The archive contains \path{frontier_experiments.py}.  It performs the
following deterministic checks at 80 decimal digits.

\begin{enumerate}
\item Coefficientwise verification of
$(I-T)^rL_r=T^{r-1}q/(1-q)$ through rank $8$ and degree $4096$.
\item Block multiplicativity of $\varepsilon_r$ through rank $8$ and large
finite blocks, plus a finite search excluding small periods.
\item Exact rational moments from Bernoulli numbers and the Bell recurrence.
\item $PF_\infty$ ultra-log-concavity ratios for ranks $1$ through $6$.
\item Numerical zero jets at $(r,Q)=(2,4),(3,8),(4,16)$.
\item Grouped Poisson-binomial means at $\theta=3^r$, with rigorous analytic
error bounds.
\item Standardized gamma cumulants through rank $64$.
\item Radial growth of $L_r$ toward an eighth root of unity.
\end{enumerate}

The exact command is
\begin{lstlisting}[style=code]
python frontier_experiments.py --output-dir numerical_output --dps 80
\end{lstlisting}
It creates
\path{numerical_results.txt}, \path{moment_ulc.csv},
\path{poisson_limit.csv}, \path{gamma_clt.csv}, and
\path{lambert_radial.csv}.

\subsection{Interpretive limits}

The finite sign search does not prove nonperiodicity; \cref{thm:nonperiodic}
does.  The radial table does not prove a natural boundary; density of the
proved dyadic radial singularities does.  The small numerical values of lower
zero derivatives reflect truncation and finite precision; their exact
vanishing follows from the multiplicity theorem.

\section{Conclusion}

The repository's existing work already connects the Fabius and Rvachev
functions to probability, exact dyadic arithmetic, infinite sinc products,
Lambert asymptotics, Legendre expansions, Bell--Bernoulli moment calculus,
$q$-binomial identities, and a higher Pascal hierarchy.  The most productive
way to advance that corpus was therefore not to derive another version of the
same endpoint or zero-count formula, but to ask what the established divisor
and automatic sign structure force in neighboring theories.

That strategy produced five coherent advances.

First, the spectral signs are not merely automatic: their generating
functions are hypertranscendental and analytically blocked by the unit circle.
Second, the zero multiplicities have their own Lambert--Mahler hierarchy with
the same strong analytic obstruction.  Third, exact lattice integration grows
from linear order in the classical rank to binomial order in the Pascal rank.
Fourth, Wick rotation promotes the entire zero divisor to total positivity at
every rank.  Fifth, the same product is a probability generating function and
a Laplace transform, yielding two new spectral probability models with exact
branching laws and rank limit theorems.

These results suggest a broader principle: the dyadic zero divisor is not just
a bookkeeping device for an infinite product.  Under different transforms it
becomes an automatic character, a Lambert $q$-series, a Strang--Fix exactness
order, a P\'olya-frequency sequence, a Poisson-binomial cloud, and a Thorin
measure.  Future progress is likely to come from moving deliberately between
those representations rather than treating them as separate topics.

\appendix

\section{Auxiliary inequalities for the Poisson phase diagram}

For reference, the elementary estimates used in
\cref{thm:PB-phase} can be organized by Bernoulli power sums.  With
$p_n=\theta/(n^2+\theta)$ and multiplicities understood,
\begin{equation}\label{eq:power-bound-app}
 S_k(\theta)=\sum\mu_r(n)p_n^k\leq\theta^k\ZetaR_r(2k).
\end{equation}
Moreover, for $k\geq1$,
\begin{equation}\label{eq:power-error-app}
 0\leq\theta^k\ZetaR_r(2k)-S_k(\theta)
 \leq k\theta^{k+1}\ZetaR_r(2k+2),
\end{equation}
because $1-(1+x)^{-k}\leq kx$ for $x\geq0$.  At a critical scaling
$\theta_r=ca^r$, the leading upper bound is
\begin{equation}\label{eq:power-scale-app}
 \theta_r^k\ZetaR_r(2k)
 =2^{2k}c^k\zeta(2k)
 \left(\frac{a^k}{4^k-1}\right)^r.
\end{equation}
This one formula explains the sequence of factorial-cumulant critical bases
$(4^k-1)^{1/k}$.

\section{Exact rational moment tables}

The following values were generated symbolically from
\eqref{eq:moment-Bell}; no floating-point reconstruction was used.

\begin{center}
\begin{tabular}{@{}rlll@{}}
\toprule
$r$ & $m_{r,0}$ & $m_{r,2}$ & $m_{r,4}$\\
\midrule
1 & $1$ & $1/9$ & $19/675$\\
2 & $1$ & $1/27$ & $107/30375$\\
3 & $1$ & $1/81$ & $571/1366875$\\
4 & $1$ & $1/243$ & $2963/61509375$\\
\bottomrule
\end{tabular}
\end{center}

\begin{center}
\small
\begin{tabular}{@{}rll@{}}
\toprule
$r$ & $m_{r,6}$ & $m_{r,8}$\\
\midrule
1 & $583/59535$ & $132809/32531625$\\
2 & $27953/56260575$ & $698600719/7839308334375$\\
3 & $1165663/53166243375$ & $2864068312649/1889077325876015625$\\
4 & $45163673/50242099989375$ & $10285851106584079/455220408602972865234375$\\
\bottomrule
\end{tabular}
\end{center}

\section{Repository and literature references}

\begin{thebibliography}{99}

\bibitem{repo-primary}
V. Reshetnikov,
\emph{The Fabius Function and Rvachev's Up-Function: Exact arithmetic,
probability, Fourier analysis, and corrected sharp asymptotics},
\texttt{ProveIt},
\url{https://github.com/VladimirReshetnikov/ProveIt/blob/main/Analysis/FabiusFunction/docs/Fabius_Function_and_Rvachev_Up/Fabius_Function_and_Rvachev_Up.tex},
repository state retrieved 30 August 2026.

\bibitem{repo-frontier}
V. Reshetnikov,
\emph{Semi-formalized research frontiers},
\url{https://github.com/VladimirReshetnikov/ProveIt/blob/main/Analysis/FabiusFunction/docs/semi-formalized-research-frontiers/semi-formalized-research-frontiers.tex},
retrieved 30 August 2026.

\bibitem{repo-spectra}
V. Reshetnikov,
\emph{Spectra and Arithmetic Frontiers},
\url{https://github.com/VladimirReshetnikov/ProveIt/blob/main/Analysis/FabiusFunction/docs/semi-formalized-research-frontiers/drafts/spectra-and-arithmetic/Spectra_and_Arithmetic_Frontiers/Spectra_and_Arithmetic_Frontiers.tex},
retrieved 30 August 2026.

\bibitem{repo-q}
V. Reshetnikov,
\emph{Exponent-Sequence and $q$-Series Frontiers},
\url{https://github.com/VladimirReshetnikov/ProveIt/blob/main/Analysis/FabiusFunction/docs/semi-formalized-research-frontiers/drafts/exponents-and-q-series/Exponents_and_q_Series_Frontiers/Exponents_and_q_Series_Frontiers.tex},
retrieved 30 August 2026.

\bibitem{repo-comb}
V. Reshetnikov,
\emph{Dyadic-Comb Frontiers for the Fabius--Rvachev System},
\url{https://github.com/VladimirReshetnikov/ProveIt/blob/main/Analysis/FabiusFunction/docs/semi-formalized-research-frontiers/drafts/inverse-and-sampling/Dyadic_Comb_Frontiers/Dyadic_Comb_Frontiers.tex},
retrieved 30 August 2026.

\bibitem{repo-compilations}
V. Reshetnikov,
\emph{Frontier Compilations},
\url{https://github.com/VladimirReshetnikov/ProveIt/blob/main/Analysis/FabiusFunction/docs/semi-formalized-research-frontiers/drafts/frontier-compilations/Frontier_Compilations/Frontier_Compilations.tex},
retrieved 30 August 2026.

\bibitem{repo-manifest}
V. Reshetnikov,
\emph{Draft manifest},
\url{https://github.com/VladimirReshetnikov/ProveIt/blob/main/Analysis/FabiusFunction/docs/semi-formalized-research-frontiers/drafts/MANIFEST.md},
retrieved 30 August 2026.

\bibitem{adh}
B. Adamczewski, T. Dreyfus, and C. Hardouin,
\emph{Hypertranscendence and linear difference equations},
J. Amer. Math. Soc. \textbf{34} (2021), no. 2, 475--503.

\bibitem{bcr}
J. P. Bell, M. Coons, and E. Rowland,
\emph{The rational--transcendental dichotomy of Mahler functions},
J. Integer Seq. \textbf{16} (2013), Article 13.2.10, 11 pp.

\bibitem{aesw}
M. Aissen, A. Edrei, I. J. Schoenberg, and A. Whitney,
\emph{On the generating functions of totally positive sequences},
Proc. Natl. Acad. Sci. USA \textbf{37} (1951), 303--307.
See also M. Aissen, I. J. Schoenberg, and A. M. Whitney,
\emph{On the generating functions of totally positive sequences I},
J. Analyse Math. \textbf{2} (1952), 93--103.

\bibitem{lecam}
L. Le Cam,
\emph{An approximation theorem for the Poisson binomial distribution},
Pacific J. Math. \textbf{10} (1960), 1181--1197.

\bibitem{stanley}
R. P. Stanley,
\emph{Differentiably finite power series},
European J. Combin. \textbf{1} (1980), 175--188.

\bibitem{allouche-shallit}
J.-P. Allouche and J. Shallit,
\emph{Automatic Sequences: Theory, Applications, Generalizations},
Cambridge University Press, 2003.

\bibitem{arias}
J. Arias de Reyna,
\emph{Arithmetic of the Fabius function},
arXiv:1702.06487, revised 2018.

\end{thebibliography}

\end{document}
